% =============================================================================================
% 4 Discussion + 5 Conclusion (v2, 2026-10-24 evening). Literature comparisons moved here from Results.
% =============================================================================================
\section{Discussion}
\label{sec:discussion}

\paragraph{The clinical model against prior work.}
Examination and symptoms carried the clinical signal: demographics alone reached AUPRC 0.09, adding the examination 0.29
and adding symptoms 0.39, whereas care, social and biomarker groups added nothing. The resulting elastic net (AUROC 0.81)
is well above the clinical GAM of the OAI study (0.68; Table~\ref{tab:main})~\cite{jang2023}. The difference is
plausibly the cohort rather than the learner. OAI is a knee cohort with little hip-specific examination, whereas CHECK
recorded a standardised hip examination and hip-specific questionnaires in people who presented early with hip
complaints. Restricted, painful internal rotation is part of the clinical definition of hip OA~\cite{altman1991acr}, and
in our data it carries more information than generic symptom scales. Deep tabular models did not beat the penalised
linear model at 93 events, as expected from the comparison literature~\cite{christodoulou2019}. The FT-Transformer's
contribution is a better calibration slope (1.04 against 0.84) and usable 2- and 5-year rankings.

\begin{table}[!tbp]
  \centering
  \caption{\textbf{Main-model performance against published multimodal models.} Our rows are internal estimates: means
  over five repeats of 5-fold nested CV, with 95\% CIs from a participant-clustered bootstrap of repeat 0. Lift = AUPRC
  $\div$ event rate. Literature values are from the abstracts and, where available, the full texts. \textsuperscript{\dag}Different endpoint
  (radiographic OA or KL progression); shown for scale only.}
  \label{tab:main}
  \small
  \setlength{\tabcolsep}{3.5pt}
  \resizebox{\linewidth}{!}{%
  \begin{tabular}{@{}llcccc@{}}
    \toprule
    Study / model & Cohort, endpoint & Event rate & AUPRC (95\% CI) & Lift & AUROC (95\% CI) \\
    \midrule
    Jang 2023, clinical GAM~\cite{jang2023} & OAI, THA 10 y & 2.4\% & 0.08 & 3.3 & 0.68 \\
    Jang 2023, clinical + DL measurements & OAI, THA 10 y & 2.4\% & 0.28 & 11.7 & 0.81 \\
    Musbahi 2026, DenseNet-201 + clinical~\cite{musbahi2026} & MOST, TKA & $\approx$10\% & 0.62 & 6.2 & 0.85 (ext.\ 0.79) \\
    Gielis 2020, clinical + Shape-Score~\cite{gielis2020}\textsuperscript{\dag} & CHECK, rHOA/THA 8 y & 17.7\% & -- & -- & 0.864 (0.833--0.895) \\
    van den Berg 2026, hybrid~\cite{vandenberg2026}\textsuperscript{\dag} & 7 cohorts, rHOA 4--8 y & 5.1\% & -- & -- & 0.81 (LOCO 0.64) \\
    Tiulpin 2019, CNN + clinical~\cite{tiulpin2019}\textsuperscript{\dag} & OAI$\to$MOST, knee KL & 38\% & 0.68 (0.66--0.70) & 1.8 & 0.79 (0.78--0.81) \\
    \midrule
    \textbf{M0} clinical elastic net & CHECK, THA 10 y & 5.5\% & 0.366 (0.281--0.487) & 6.7 & 0.807 (0.771--0.868) \\
    M0 + frozen DenseNet PCA-32 (v1) & CHECK, THA 10 y & 5.5\% & 0.386 (0.301--0.499) & 7.1 & 0.820 (0.777--0.873) \\
    \textbf{M0 + pixel logit (final)} & CHECK, THA 10 y & 5.5\% & \textbf{0.552} (0.472--0.685) & \textbf{10.1} & \textbf{0.892} (0.861--0.936) \\
    \bottomrule
  \end{tabular}}
\end{table}

\paragraph{Transport across centres.}
Discrimination transported across pseudo-centres: under leave-one-centre-out validation M0 kept its AUPRC (0.374 against
0.377 internally) and lost 0.02 AUROC, and the image head lost 0.02 AUROC. Published hip models fare worse under the
same design. A seven-cohort model fell from AUROC 0.81 to 0.64 under leave-one-cohort-out validation~\cite{vandenberg2026},
and a knee surgical model fell from 0.85 to 0.79 externally~\cite{musbahi2026}. Clinical-only knee models are known to
transport well~\cite{castagno2025}, and our result extends this to the hip with a standardised examination. The
comparison flatters us, however. Our pseudo-centres share country, era, protocol and readers, so the shift between them
is milder than between cohorts. Calibration did not transport (intercepts $-1.6$ to $+1.4$; painless hips
over-predicted), and image discrimination was sensitive to film format. Local recalibration is therefore a precondition
for use, as for most prediction models~\cite{steyerberg2016validation}.

\paragraph{The image numbers in a rare-event setting.}
At a 5.5\% event rate, an image-only AUROC of 0.87 and a fusion gain of $+0.20$ AUPRC (0.37 $\to$ 0.55) must be read on
the precision scale. The highest attainable lift is the inverse of the event rate: 18 here and 42 in OAI. Our fused lift
of 10.1 thus uses about half of the available range, and Jang's 11.7 about a quarter of its range. The image-only head
(lift 9.0) matches Jang's measurement model (9.2) at twice the event rate (Table~\ref{tab:image}). The more
important question is where the gain comes from. Radiographic severity is a core indication for THA: surgeons offer
replacement when symptoms persist \emph{and} the film shows advanced joint damage~\cite{gademan2016indication}. The
baseline film therefore carries part of the decision that defines the outcome, and the gain is concentrated in hips
replaced within 1--3 years (AUROC 0.90--0.98 for every model). In our reading, the improvement reflects what the image
contains rather than an advantage of attention as a mechanism. The architecture mattered only for extraction. The same
pixels gave AUPRC 0.12 through frozen pooled features and 0.49 through the spatial head.

\begin{table}[!tbp]
  \centering
  \caption{\textbf{Image-only performance.} Same conventions as Table~\ref{tab:main}.
  \textsuperscript{\dag}Different endpoint (shown for scale only).}
  \label{tab:image}
  \small
  \setlength{\tabcolsep}{3.5pt}
  \resizebox{\linewidth}{!}{%
  \begin{tabular}{@{}llcccc@{}}
    \toprule
    Study / model & Cohort, endpoint & Event rate & AUPRC (95\% CI) & Lift & AUROC (95\% CI) \\
    \midrule
    Agricola 2013, shape modes (apparent)~\cite{agricola2013} & CHECK, THA 5 y & 2.5\% & -- & -- & 0.81 \\
    Jang 2023, 26 DL measurements~\cite{jang2023} & OAI, THA 10 y & 2.4\% & 0.22 & 9.2 & 0.77 \\
    van den Berg 2026, DL features~\cite{vandenberg2026}\textsuperscript{\dag} & 7 cohorts, rHOA 4--8 y & 5.1\% & -- & -- & 0.81 \\
    Tiulpin 2019, CNN~\cite{tiulpin2019}\textsuperscript{\dag} & OAI$\to$MOST, knee KL & 38\% & 0.68 (0.66--0.70) & 1.8 & 0.79 (0.77--0.80) \\
    \midrule
    Frozen DenseNet-121, GAP + PCA-32 & CHECK, THA 10 y & 5.5\% & 0.122 (0.098--0.222) & 2.2 & 0.687 (0.649--0.758) \\
    Linear head on GAP features & CHECK, THA 10 y & 5.5\% & 0.228 (0.176--0.370) & 4.2 & 0.769 (0.723--0.829) \\
    \textbf{Residual head + OARSI targets} & CHECK, THA 10 y & 5.5\% & \textbf{0.491} (0.403--0.624) & \textbf{9.0} & \textbf{0.866} (0.834--0.917) \\
    \bottomrule
  \end{tabular}}
\end{table}

\paragraph{Why the image head works, and what it may be reading.}
Several observations support a genuine pixel signal. Most of the gain came from the spatial residual stage (linear
$\to$ residual: $+0.27$ AUPRC). Attention concentrated on the superolateral joint space, the typical site of early
narrowing in hip OA. The head recovered OARSI items better than frozen features did (KL $\rho$ 0.31 $\to$ 0.45). A head
trained without the reading targets fused equally well (AUPRC 0.527 against 0.552), and without THA in years 1--2 the
fusion gain remained $+0.20$. Other observations argue for caution. False positives were KL 1--2 hips with attention in
the same place. This suggests the head encodes radiographic stage, which is partly the surgeon's decision rule, rather
than biology that predicts progression. In missed hips the attention drifted away from the joint, but confining it to
the joint did not recover them. The drift is therefore a symptom of finding no signal, not the cause of the miss.
Crop embeddings still identified the centre and the format. Attention maps are descriptive, not attributions. With observational data we cannot separate these explanations causally.

\paragraph{Why token fusion failed.}
Fusing the image as one FT-Transformer token lowered AUPRC by 0.13 relative to M0. The likely cause is the input to
the token, not the transformer. M4 received the frozen embedding after global pooling and PCA to 32 dimensions. That
signal was weak on its own (AUPRC 0.12), because pooling discards the location of millimetre-scale joint-space changes.
The token's input LayerNorm then normalised each hip's 32-vector, which removes its magnitude. With about 60 events
per training fold, the attention layers had to learn a cross-modal interaction from a weak, noisy token. A follow-up
run supports this reading. Feeding the cross-fitted pixel logit as the token (with the LayerNorm bypassed) raised the
FT-Transformer from 0.323 to 0.471 AUPRC ($+0.147$, 0.038--0.250; AUROC 0.848, calibration slope 0.99). Its 10-year
discrete-time risk reached AUROC 0.83. It remained below the elastic-net fusion ($-0.095$, $-0.150$ to $-0.043$). The
transformer is therefore not useless, but at this sample size an extra layer of interactions costs more than it gains.

\paragraph{Clinical implications.}
The model gives a first-visit risk score from routine data and the raw film, without a radiologist's grading. Its
intended uses are setting the intensity of conservative care and the review interval, timing referral, and enriching
joint-preservation trials with high-risk hips. It adds little for hips already on a surgical pathway, where THA follows
within two years and any model performs well. Before use elsewhere it needs local recalibration, especially for hips
without pain.

\paragraph{Limitation 1: data scale against model capacity.}
Deep networks learn their representations from very large data sets: ImageNet has 1.2 million labelled images. We
have 1,705 hips with known outcome, and for a binary outcome the effective sample size is closer to the 93 events. On
the course's rule of thumb (Lecture 2, slide 107), our 1,968 films sit in the regime where fine-tuning works, but the
93 positives sit in the ``marginal'' range of high dozens to low hundreds. We therefore never trained a deep network
end to end. The image backbone was pretrained on ImageNet and frozen. Only a
0.21\,M-parameter head was trained on pixels, and the FT-Transformer had 34k--56k parameters. Even so, the
FT-Transformer did not beat a penalised linear model, its repeats were unstable until seeds were averaged, and the
fusion model's calibration slope of 0.79 shows residual overfitting. The event count, not the architecture, limits
what this study can show about deep models. It does not explain the AUPRC level, which is bounded by the 5.5\% event
rate, as discussed above.

\paragraph{Further limitations of the statistical design.}
Three further issues could not be resolved within this study. First, events per parameter: 93 events against 76
features plus missing indicators is about one event per candidate parameter, below current
guidance~\cite{riley2020samplesize}. Penalisation limits overfitting but does not remove it. Second, imputation: fold-wise MICE is a single imputation under a missing-at-random assumption, so it understates
uncertainty. It also excludes the outcome, which biases imputed predictors' effects towards zero~\cite{moons2006imputation,
vanbuuren2011mice}. Collinearity is a minor concern here. BayesianRidge shrinks correlated predictors, and correlated
predictors make imputation more accurate. Because only 2--5\% of values were imputed, the impact is small. Third,
censoring: excluding the 299 hips censored before year 10 conditions on follow-up, and death as a competing risk was
not recorded.

\paragraph{Limitations of the image supervision and the stacking.}
The OARSI targets may reduce generalisation. They encode the conventions of CHECK's readers and a paired reading with
later films~\cite{macri2022}, so the head may learn features that do not transfer to cohorts read differently. World
COACH faced the same issue: its best image features were optimised with follow-up grades~\cite{vandenberg2026}. Our
no-target head matched discrimination internally, but its transport was not tested externally. The cross-fitting
itself does not leak test labels. Each training hip's logit comes from a head that never saw that hip, test hips are
scored by a head trained only on training hips, and test labels are used only for evaluation. Two weaker effects
remain. Training logits come from heads fitted on two thirds of the data, so their scale differs slightly from the test
logits. During inner tuning, inner-validation labels have shaped other training hips' logits. Both effects can bias
the choice of hyper-parameters, but not the outer-fold estimate.

\paragraph{Limitations of late fusion.}
The final model fuses late: the image head is trained on pixels alone, and its logit enters the elastic net as one
feature. Because the logit is fitted jointly with the clinical features, its coefficient is already adjusted for
their linear correlation. Late fusion still has two costs. The head never sees the clinical context, so part of what
it learns duplicates the examination. And a single linear term cannot express interactions, for example a pixel
signal that matters more in painful hips. Joint fusion is the alternative. Here it lost at this sample size, both as a
frozen-embedding token and FiLM (M4) and as a pixel-logit token (AUPRC 0.471 against 0.552). Training the head on the
residual of the clinical model, or adding pre-specified interaction terms, are cheaper middle grounds for future work.

\paragraph{Limitations of the cohort.}
CHECK enrolled Dutch adults aged 45--65 with early complaints between 2002 and 2005. Of these, 79\% were women and 97\%
were white, and more than half had not seen any health professional for the complaint in the year before baseline.
Referral passes through general practitioners, and both the threshold for surgery and access to it are specific to
that health system and era. The model may therefore not hold for older patients, for populations with different hip
morphology or OA prevalence (for example, East Asian cohorts with more dysplasia-related OA), or where access to
surgery differs. Validation was internal only. The pseudo-centres were reconstructed from device serials, two of them
impurely. Label-free steps (orientation, locator, SimCLR, AdaBN) were fitted on all films.

\paragraph{Future work.}
The first priority is external validation, in OAI hip films and in the World COACH consortium. Both would test format
and population transport, and both would show whether local recalibration is enough. On the imaging side, end-to-end
fine-tuning with an L2-SP penalty would let the backbone adapt to radiographs without forgetting its pretraining.
Higher resolution ($28\times28$ maps or $512^2$ crops) would target cysts and head flattening, which the current map
resolves poorly. Explicit geometric measurements (lateral centre-edge angle, $\alpha$ angle, minimum joint-space width)
would make the pixel score interpretable, and a contralateral-hip reference could separate a person's anatomy from
disease. Format harmonisation should reduce the TIFF/DICOM gap. Finally, decision thresholds should be set with
clinicians through decision-curve analysis rather than chosen by us.

\section{Conclusion}
\label{sec:conclusion}
We asked whether a first-consultation hip can be scored for ten-year replacement risk from routine information alone. In
the CHECK study, a penalised logistic model on the standardised hip examination and symptoms already discriminates well
(AUROC 0.81), and deep tabular models do not improve on it at 93 events. Frozen ImageNet features of the baseline film
add nothing. A small residual head that keeps the spatial map of a frozen DenseNet, trained with OARSI items as
targets only, extracts a strong pixel signal. Fused as one cross-fitted logit, it raises AUPRC from 0.37 to 0.55 in
internal validation, and the gain survives the exclusion of early events. Much of this signal reflects radiographic
stage, which is itself a criterion for surgery. Discrimination transported across pseudo-centres, but calibration did
not. The model is therefore a candidate for external validation and local recalibration, not a finished tool.
